%% ============================================================
%%  TEMPLATE TUGAS AKHIR ITS  (Program Sarjana)
%%  Kelas TA-ITS.cls (Matematika ITS), disusun ulang dengan
%%  struktur folder yang sama seperti template tesis.
%%
%%  File utama -- kompilasi file ini dengan pdfLaTeX + Biber:
%%    latexmk -pdf TA.tex
%%  Petunjuk lengkap: lihat Panduan-Template.tex / Panduan-Template.pdf
%% ============================================================

%% File kelas & paket template disimpan di folder Format/
\makeatletter
\def\input@path{{Format/}}
\makeatother
\documentclass{TA-ITS}

\theoremstyle{definition}
\newtheorem{definisi}{Definisi}[section]
\newtheorem{teorema}{Teorema}[section]
\newtheorem{lema}{Lema}[section]
\newtheorem{akibat}{Akibat}[section]
\newtheorem{contoh}{Contoh}[section]
\newtheorem{proposisi}{Proposisi}[section]
\newtheorem{catatan}{Catatan}[section]
\renewcommand{\proofname}{Bukti}

\providecommand{\definisiautorefname}{Definisi}
\providecommand{\proposisiautorefname}{Proposisi}
\providecommand{\lemaautorefname}{Lema}
\providecommand{\teoremaautorefname}{Teorema}
\providecommand{\akibatautorefname}{Akibat}
\providecommand{\contohautorefname}{Contoh}
\providecommand{\catatanautorefname}{Catatan}
\providecommand{\algorithmautorefname}{Algoritma}

\setlength{\parindent}{0.75cm} % Indentasi 0.75cm pada tiap paragraf

%% Nomor halaman: ganjil di kanan bawah, genap di kiri bawah
\usepackage{fancyhdr}
\pagestyle{fancy}
\renewcommand{\headrulewidth}{0pt}
\fancyhf{}
\fancyfoot[R]{\ifthenelse{\isodd{\value{page}}}{\thepage}{}}
\fancyfoot[L]{\ifthenelse{\isodd{\value{page}}}{}{\thepage}}

\usepackage[labelsep=quad]{caption}
\captionsetup[table]{skip=5pt}

\usepackage{multirow}
\usepackage{makecell}
\usepackage{longtable}
\usepackage{physics}
\usepackage{ragged2e}
\usepackage[mathscr]{eucal}
\usepackage[table]{xcolor}

%% Penulisan kode program (lihat contoh di Bagian_Akhir/lampiran.tex)
\usepackage{listings}
\definecolor{codegreen}{rgb}{0,0.6,0}
\definecolor{codeblack}{rgb}{0,0,0}
\definecolor{codepurple}{rgb}{0.58,0,0.82}
\lstdefinestyle{mystyle}{
    commentstyle=\color{codegreen},
    keywordstyle=\color{magenta},
    numberstyle=\tiny\color{codeblack},
    stringstyle=\color{codepurple},
    basicstyle=\ttfamily\footnotesize,
    breaklines=true,
    captionpos=b,
    keepspaces=true,
    numbers=left,
    numbersep=5pt,
    showstringspaces=false,
    tabsize=2
}

%% Diagram alir (gaya didefinisikan di Bab/Bab_3/flowchart.tex)
\usetikzlibrary{shapes.geometric,arrows}

\newcommand{\R}{\mathbb{R}}
\newcommand{\Z}{\mathbb{Z}}
\newcommand{\N}{\mathbb{N}}
\newcommand{\C}{\mathbb{C}}
\newcommand{\Q}{\mathbb{Q}}

%% Tambahkan paket / perintah buatan sendiri di sini
%% contoh: \newcommand{\Fix}{\operatorname{Fix}}

\addbibresource{Pustaka/daftarPustaka.bib}

\begin{document}

%% ============================================================
%%  DATA TUGAS AKHIR  (wajib diisi)
%% ============================================================
\Judul{Judul Tugas Akhir dalam Bahasa Indonesia}
\JudulEng{Final Project Title in English}
\Nama{NAMA LENGKAP MAHASISWA}                % sampul (otomatis kapital)
\NamaKecil{Nama Lengkap Mahasiswa}           % lembar lain (huruf biasa)
\NRP{50022XXXXX}
\KodeTA{SM234801}                            % kode mata kuliah TA (sampul)

\ProgramStudi{Matematika}                    % prodi (bhs. Indonesia)
\ProgramStudiEng{Mathematics}                % prodi (bhs. Inggris)
\Gelar{Sarjana Matematika}                   % gelar pada lembar pengesahan
\GelarEng{Bachelor of Mathematics}
\Departemen{Matematika}
\Department{Mathematics}
\BidangStudi{Analisis dan Aljabar}
\Fakultas{Sains dan Analitika Data}{FSAD}    % Fakultas {nama}{singkatan}
\Faculty{Science and Data Analytics}
\SingkatanFakultas{FSAD}                     % dipakai di identitas abstrak
\SingkatanFakultasEng{SCIENTICS}

\Bulan{Juli}                                 % bulan pengesahan (bhs. Inggris otomatis)
\Tahun{20XX}
\TanggalDisetujui{DD Bulan 20XX}             % tanggal pada pernyataan orisinalitas
\ApprovalDate{DD Month 20XX}                 % idem, versi bhs. Inggris

\Pembimbing{Nama Pembimbing 1, gelar}        % Pembimbing 1
{}                                           % Pembimbing 2 (kosongkan bila tidak ada)
{}                                           % Pembimbing 3 (kosongkan bila tidak ada)

\NIPPembimbing{XXXXXXXX XXXXXX X XXX}
{}
{}

\Penguji{Nama Penguji 1, gelar}
{Nama Penguji 2, gelar}
{Nama Penguji 3, gelar}

\NIPPenguji{XXXXXXXX XXXXXX X XXX}
{XXXXXXXX XXXXXX X XXX}
{XXXXXXXX XXXXXX X XXX}

\Kadep{Nama Kepala Departemen, gelar}
\NIPKadep{XXXXXXXX XXXXXX X XXX}
\Kaprodi{Nama Kepala Program Studi, gelar}   % hanya dipakai \LembarPengesahanProposal
\NIPKaprodi{XXXXXXXX XXXXXX X XXX}

\allowdisplaybreaks

%% ============================================================
%%  BAGIAN AWAL
%% ============================================================
\Cover                 % sampul luar (biru ITS)
\BagianAwal
\LembarJudul           % halaman judul bahasa Indonesia
\TitlePage             % halaman judul bahasa Inggris
\LembarPengesahan
\LembarPengesahanEng
% \LembarPengesahanProposal  % pakai ini (ganti dua baris di atas) untuk proposal TA
\LembarOrisinalitas
\LembarOrisinalitasEng

\LembarAbstrak{Bagian_Awal/abstrak}
\LembarAbstract{Bagian_Awal/abstract}
\KataPengantar{Bagian_Awal/katapengantar}

\DaftarIsi\raggedbottom
\DaftarGambar
\DaftarTabel
\DaftarSimbol{Bagian_Awal/daftarsimbol}

%% ============================================================
%%  BAGIAN INTI
%% ============================================================
\BagianInti

\chapter{PENDAHULUAN}
Pada bab ini dijelaskan tentang latar belakang, rumusan masalah, batasan masalah, tujuan penelitian, dan manfaat penelitian.

\section{Latar Belakang}
Uraikan konteks penelitian, untuk apa penelitian ini dilakukan, dan hal apa yang mengarahkan penelitian ini. Latar belakang memuat teori utama yang relevan dan terbaru \parencite{contoh2024artikel} yang dipadukan sehingga mengerucut pada suatu persoalan yang kemudian disusun dalam bentuk rumusan masalah. Sitasi naratif ditulis dengan \texttt{\textbackslash textcite}, misalnya \textcite{penulis2015buku}.

\section{Rumusan Masalah}
Rumusan masalah dalam penelitian ini adalah sebagai berikut:
\begin{enumerate}
  \item Rumusan masalah pertama?
  \item Rumusan masalah kedua?
\end{enumerate}

\section{Batasan Masalah}
Batasan masalah dalam penelitian ini adalah sebagai berikut:
\begin{enumerate}
  \item Batasan pertama.
  \item Batasan kedua.
\end{enumerate}

\section{Tujuan Penelitian}
Berdasarkan rumusan masalah, tujuan dari penelitian ini adalah sebagai berikut:
\begin{enumerate}
  \item Tujuan yang menjawab rumusan masalah pertama.
  \item Tujuan yang menjawab rumusan masalah kedua.
\end{enumerate}

\section{Manfaat Penelitian}
Manfaat yang diharapkan dari penelitian ini adalah sebagai berikut:
\begin{enumerate}
  \item Secara teoritis, \dots
  \item Secara praktis, \dots
  \item Sebagai bahan referensi untuk penelitian selanjutnya.
\end{enumerate}

\chapter{TINJAUAN PUSTAKA}
Pada bab ini dipaparkan penelitian terdahulu yang berkaitan dengan penelitian ini serta definisi, lema, teorema, dan contoh yang menjadi dasar teori.

\section{Hasil Penelitian Terdahulu}
Beberapa penelitian terdahulu yang berkaitan dengan penelitian ini dirangkum pada \ref{tab:penelitian_terdahulu}.

\begin{table}[H]
  \centering
  \caption{Penelitian terdahulu}
  \label{tab:penelitian_terdahulu}
  \begin{tabular}{|c|m{4cm}|c|m{8cm}|}
    \hline
    \textbf{No} & \textbf{Peneliti}                           & \textbf{Tahun}                 & \textbf{Hasil Penelitian}            \\ \hline
    1           & \citeauthorfullname{penulis2015buku}        & \citeyear{penulis2015buku}     & Ringkasan hasil penelitian pertama.  \\ \hline
    2           & \citeauthorfullname{contoh2024artikel}      & \citeyear{contoh2024artikel}   & Ringkasan hasil penelitian kedua.    \\ \hline
    3           & \citeauthorfullname{mahasiswa2023skripsi}   & \citeyear{mahasiswa2023skripsi} & Ringkasan hasil penelitian ketiga.  \\ \hline
  \end{tabular}
\end{table}

\section{Dasar Teori}
Pada bagian ini diberikan definisi, contoh, lema, dan teorema yang digunakan dalam penelitian.

\begin{definisi}[\cite{penulis2015buku}]\label{def:metrik}
  Diberikan himpunan tak kosong $X$. Fungsi $d:X\times X\to\R$ disebut metrik pada $X$ jika untuk setiap $x,y,z\in X$ berlaku
  \begin{enumerate}
    \item[(a)] $d(x,y)\geq 0$ dan $d(x,y)=0$ jika dan hanya jika $x=y$,
    \item[(b)] $d(x,y)=d(y,x)$,
    \item[(c)] $d(x,z)\leq d(x,y)+d(y,z)$.
  \end{enumerate}
\end{definisi}

\begin{contoh}\label{con:metrik}
  Fungsi $d(x,y)=|x-y|$ untuk setiap $x,y\in\R$ merupakan metrik pada $\R$.
\end{contoh}

\subsection{Contoh Subbab}
\begin{lema}\label{lem:contoh}
  Tuliskan pernyataan lema di sini.
\end{lema}
\begin{proof}
  Tuliskan bukti lema di sini.
\end{proof}

\begin{teorema}[\cite{contoh2024artikel}]\label{thm:contoh}
  Tuliskan pernyataan teorema di sini, misalnya
  \begin{equation}\label{eq:contoh}
    \norm{x+y}\leq\norm{x}+\norm{y}.
  \end{equation}
\end{teorema}

Berdasarkan \autoref{def:metrik}, \autoref{con:metrik}, dan \autoref{thm:contoh}, khususnya Persamaan \eqref{eq:contoh}, \dots


\chapter{METODOLOGI}
Pada bab ini dijelaskan langkah-langkah pelaksanaan penelitian Tugas Akhir beserta diagram alir penelitian.

\section{Tahapan Penelitian}
\begin{itemize}
    \item[(a)] Studi Literatur\\
    Tahap ini dilakukan dengan mempelajari literatur yang berkaitan dengan topik penelitian, baik berupa buku, jurnal, artikel, maupun sumber terpercaya lainnya.
    \item[(b)] Tahap Kedua\\
    Jelaskan kegiatan pada tahap ini.
    \item[(c)] Tahap Ketiga\\
    Jelaskan kegiatan pada tahap ini.
    \item[(d)] Penarikan Kesimpulan\\
    Kesimpulan disusun berdasarkan hasil pada tahap-tahap sebelumnya untuk menjawab rumusan masalah.
    \item[(e)] Penulisan Tugas Akhir\\
    Tahap akhir dari penelitian ini adalah menyusun laporan Tugas Akhir berdasarkan hasil penelitian yang telah dilakukan.
\end{itemize}

\section{Diagram Alir Penelitian}
Diagram alir penelitian disajikan pada \ref{fig:diagram_alir} sebagai berikut.
\begin{figure}[H]
    \centering
    \tikzstyle{startstop} = [ellipse, minimum width=5.5cm, minimum height=0.5cm, text centered, draw=black]
\tikzstyle{process} = [rectangle, minimum height=0.5cm, text centered, text width=8.5cm, draw=black]
\tikzstyle{arrow} = [thick,->,>=stealth]
\begin{tikzpicture}[node distance=1.3cm]
\node (start) [startstop] {Mulai};
\node (pro1) [process, below of=start] {Studi Literatur};
\node (pro2) [process, below of=pro1] {Tahap Kedua};
\node (pro3) [process, below of=pro2] {Tahap Ketiga};
\node (pro4) [process, below of=pro3] {Penarikan Kesimpulan};
\node (pro5) [process, below of=pro4] {Penulisan Tugas Akhir};
\node (stop) [startstop, below of=pro5] {Selesai};
\draw [arrow] (start) -- (pro1);
\draw [arrow] (pro1) -- (pro2);
\draw [arrow] (pro2) -- (pro3);
\draw [arrow] (pro3) -- (pro4);
\draw [arrow] (pro4) -- (pro5);
\draw [arrow] (pro5) -- (stop);
\end{tikzpicture}

    \caption{Diagram alir penelitian}
    \label{fig:diagram_alir}
\end{figure}

\chapter{PEMBAHASAN}
Pada bab ini dibahas hasil penelitian yang menjawab rumusan masalah.

\section{Hasil Teoritis}
Tuliskan hasil utama penelitian beserta buktinya.

\begin{proposisi}\label{prop:utama}
  Tuliskan pernyataan proposisi di sini.
\end{proposisi}
\begin{proof}
  Tuliskan bukti proposisi di sini.
\end{proof}

\begin{akibat}
  Akibat langsung dari \autoref{prop:utama}.
\end{akibat}

\begin{catatan}
  Catatan tambahan terkait hasil di atas.
\end{catatan}

\section{Algoritma dan Simulasi}
Langkah perhitungan dapat dituliskan sebagai algoritma, seperti pada \autoref{alg:contoh}.

\begin{algorithm}
  \caption{Contoh algoritma}\label{alg:contoh}
  \begin{algorithmic}[1]
    \Require $x_0\in\R$, toleransi $\varepsilon>0$
    \Ensure Hampiran $x^*$
    \State $k\gets 0$
    \Repeat
      \State $x_{k+1}\gets T(x_k)$
      \State $k\gets k+1$
    \Until{$|x_{k}-x_{k-1}|<\varepsilon$}
    \State \Return $x_k$
  \end{algorithmic}
\end{algorithm}

Hasil simulasi disajikan pada \ref{tab:hasil}. Judul tabel diletakkan di atas tabel, sedangkan judul gambar di bawah gambar.

\begin{table}[H]
  \centering
  \caption{Contoh tabel hasil simulasi}
  \label{tab:hasil}
  \renewcommand{\arraystretch}{1.2}
  \begin{tabular}{|c|c|c|}
    \hline
    \textbf{Iterasi} & \textbf{$x_k$} & \textbf{Galat} \\ \hline
    1 & 0,500 & 0,500 \\ \hline
    2 & 0,750 & 0,250 \\ \hline
    3 & 0,875 & 0,125 \\ \hline
  \end{tabular}
\end{table}

\begin{figure}[H]
  \centering
  \includegraphics[width=0.3\textwidth]{lambang-its}
  \caption{Contoh gambar (sumber: dokumen pribadi)}
  \label{fig:contoh}
\end{figure}


\chapter{KESIMPULAN DAN SARAN}

\section{Kesimpulan}
Berdasarkan pembahasan pada Bab IV, dapat ditarik kesimpulan sebagai berikut.
\begin{enumerate}
  \item Kesimpulan yang menjawab rumusan masalah pertama.
  \item Kesimpulan yang menjawab rumusan masalah kedua.
\end{enumerate}

\section{Saran}
Berikut merupakan saran untuk penelitian selanjutnya.
\begin{enumerate}
  \item Saran pertama.
  \item Saran kedua.
\end{enumerate}


\DaftarPustaka

%% ============================================================
%%  BAGIAN AKHIR
%% ============================================================
\DaftarLampiran{Bagian_Akhir/lampiran}
\UcapanTerimaKasih{Bagian_Akhir/terimakasih}   % opsional
\Biodata{foto}                                 % foto penulis 4x6 di Gambar/
Nama Lengkap Mahasiswa, lahir di Kota pada DD Bulan YYYY. Penulis menempuh pendidikan formal di SD \dots, SMP \dots, dan SMA \dots. Penulis diterima di Departemen Matematika, Institut Teknologi Sepuluh Nopember (ITS) Surabaya melalui jalur \dots\ pada tahun YYYY. Selama kuliah, penulis aktif dalam kegiatan \dots. Penulis dapat dihubungi melalui surel \texttt{nama@email.com}.


\end{document}
